\section*{\centering\heiti\zihao{4}摘\quad 要}
\addcontentsline{toc}{section}{摘要}

\vspace{0.5cm}

{\zihao{-4}
面向金融反诈智能守护赛题一“欺诈交易智能识别”的要求，本文档给出FP-FraudSim流式风控系统的完整技术方案与工程实现说明。系统围绕第三方支付跨时段、跨渠道交易风险防控场景，构建从仿真数据生成、欺诈模式注入、流式实时识别、风险等级输出、团伙溯源分析到人工反馈迭代的端到端闭环，目标是在交易发生时完成低延迟风险评分，并为人工研判提供可解释证据。

系统采用Kafka + Flink + Redis + FastAPI + Dashboard的工程架构。离线侧通过公开交易仿真数据规范化与场景注入，构造账户接管、钓鱼转账、洗钱拆分、卡农账户、优惠滥用、商户套现、设备团伙欺诈等七类样本；训练侧融合用户画像、设备指纹、IP环境、商户行为、短时窗口统计和图关联特征，支持LightGBM、XGBoost、CatBoost、sklearn等多模型训练与排行榜管理；在线侧由Flink消费交易流，实时计算5分钟、10分钟、1小时窗口特征，从Redis补齐画像和图风险，调用模型服务输出风险分数、风险等级、处置动作与理由码，并将结果写入risk\_results、alert\_events和late\_events等主题。

围绕赛题评分准则，系统重点实现四项能力：一是可重复的多源异构仿真环境与跨渠道交易流回放；二是Kafka + Flink微批评分链路与批量推理接口，兼顾吞吐和延迟；三是人工反馈与轻量重训机制，支持新增欺诈样本回流和模型热切换；四是基于共享设备、IP、商户、收款方的图挖掘模块，输出疑似团伙、风险得分和证据链。工程层面提供Docker Compose编排、一键启动脚本、API压测脚本、单元测试与可视化看板，便于初赛复现和现场答辩演示。

本文档的贡献在于将赛题一要求拆解为可运行系统模块，并在文档中逐项对齐完成度、创新性、实用性、成熟度和安全性评分点。系统不是单一离线分类模型，而是具备交易仿真、实时检测、动态异常识别、可解释溯源和持续迭代能力的金融反诈原型。
}

\vspace{0.6cm}

\noindent{\heiti\zihao{-4}关键词：}{\zihao{-4}金融反诈；欺诈交易识别；Kafka；Flink；图挖掘；模型迭代}
